\begin{table}[H]
\centering
\caption{Standardized Effect Sizes for Main Outcomes}
\label{tab:sde}
\begin{adjustbox}{max width=\textwidth}
\begin{tabular}{llcccccl}
\toprule
Outcome & Specification & $\hat{\beta}$ & SD($X$) & SD($Y$) & SDE & SE(SDE) & Classification \\
\midrule
Child labor rate & DiD (1910--1920) & 0.2154 & --- & 0.098 & 2.191 & 0.422 & Large positive \\
School attendance & DiD (1910--1920) & -0.1338 & --- & 0.057 & -2.366 & 0.411 & Large negative \\
SEI (1940) & OLS, continuous & 0.1041 & 2.855 & 3.189 & 0.093 & 0.129 & Moderate positive \\
Occscore (1940) & OLS, continuous & 0.3260 & 2.855 & 1.705 & 0.546 & 0.117 & Large positive \\
School attend. (1930) & OLS, continuous & 0.0037 & 2.855 & 0.037 & 0.280 & 0.181 & Large positive \\
Farm residence (1940) & OLS, continuous & -0.0046 & 2.855 & 0.135 & -0.097 & 0.092 & Moderate negative \\
\bottomrule
\end{tabular}
\end{adjustbox}
\par\vspace{0.3em}
{\footnotesize \emph{Notes:} This table reports standardized effect sizes (SDE) to facilitate cross-study comparison of treatment effect magnitudes. For binary (0/1) treatments, SDE $= \hat{\beta} / \text{SD}(Y)$ and the SD($X$) column is marked ``---''. For continuous treatments, SDE $= \hat{\beta} \times \text{SD}(X) / \text{SD}(Y)$, which gives the effect of a one-standard-deviation change in the treatment variable, measured in standard deviations of the outcome. SD($Y$) and SD($X$) are unconditional standard deviations from the summary statistics (Table~\ref{tab:summary}), before conditioning on fixed effects. 

\textbf{Research question:} Does childhood exposure to mothers' pensions (America's first welfare program, adopted 1911--1919) improve children's long-run occupational attainment? 
\textbf{Treatment:} Binary (state adopted MP by 1919) for child labor/schooling; continuous (years of MP exposure = 1920 $-$ adoption year) for adult outcomes. 
\textbf{Data:} IPUMS Multigenerational Longitudinal Panel, 1910--1940 linked censuses, state-level aggregates. 
\textbf{Method:} DiD (child labor, schooling); OLS with controls (adult outcomes). State-clustered SEs. 
\textbf{Sample:} Children aged 8--14 in 1910 (short-run) or 0--10 in 1920 (long-run), tracked across census waves.

Classification thresholds (7 categories): large negative ($< -0.15$), moderate negative ($-0.15$ to $-0.05$), small negative ($-0.05$ to $-0.005$), null ($-0.005$ to $0.005$), small positive ($0.005$ to $0.05$), moderate positive ($0.05$ to $0.15$), large positive ($> 0.15$). Classification labels refer to the magnitude of the standardized point estimate, not to statistical significance. ``Null'' denotes a near-zero effect size ($|$SDE$| < 0.005$), not a failure to reject a null hypothesis.}
\end{table}
